\documentclass[11pt]{article}
\usepackage{amsmath,amssymb,amsthm,hyperref}
\newtheorem{proposition}{Proposition}
\newtheorem{corollary}[proposition]{Corollary}
\title{A quadratic limitation of disjoint linear moment corrections}
\author{Research note, 30 September 2026}
\date{}
\begin{document}
\maketitle

\noindent\textbf{Status.} Complete elementary proof. This is a limitation
of the new construction method, not a lower bound for arbitrary fair dice.
It explains why reducing the number of nodes in each private correction
set cannot turn that method into a linear-size exceptional die.

\begin{proposition}\label{prop:barrier}
Let $m\geq2$, and let $B$ be a nonzero linear functional on the polynomials
of degree at most $m-1$. Suppose that, for each $j\in[m]$, it has a signed
atomic representation
\[
 B(a)=\sum_{q\in Q_j}\delta_{j,q}a(t_q)
 \qquad(\deg a\leq m-1),
\]
where all $t_q$ are distinct, the sets $Q_1,\ldots,Q_m$ are pairwise
disjoint, and all displayed weights are nonzero. Then
\[
 |Q_i|+|Q_j|\geq m+1\quad(i\ne j),\qquad
 \sum_{j=1}^m|Q_j|\geq\frac{m(m+1)}2.
\]
More precisely, minimizing the integer sum under these pair inequalities
requires at least
\[
 \begin{cases}
  m(m+1)/2,&m\text{ odd},\\
  m(m+2)/2-1,&m\text{ even}.
 \end{cases}
\]
\end{proposition}
\begin{proof}
Subtract the representations for $i$ and $j$. This gives a nonzero signed
measure with support $Q_i\cup Q_j$ whose moments of orders
$0,\ldots,m-1$ vanish. A signed measure supported on $s\leq m$ distinct
points and with these vanishing moments must be zero: apply the first $s$
moment equations, whose coefficient matrix is an invertible Vandermonde
matrix. Since the supports are disjoint and all weights are nonzero,
the difference measure is nonzero. Thus $|Q_i|+|Q_j|\geq m+1$.
Summing over unordered pairs gives
\[
 (m-1)\sum_j|Q_j|\geq\binom m2(m+1).
\]
For the sharper integer statement, write $a_1\leq\cdots\leq a_m$ for
the support sizes. Only $a_1+a_2\geq m+1$ matters for minimization.
If $m=2r+1$, taking all $a_j=r+1$ attains the displayed bound.
If $m=2r$, either $a_1\geq r+1$, giving total at least $m(r+1)$,
or $a_1\leq r$ and $a_j\geq2r+1-a_1$ for $j\geq2$.
The latter sum is minimized at $a_1=r$, and equals
$r+(m-1)(r+1)=m(m+2)/2-1$.
\end{proof}

\paragraph{Application to the fair-dice construction.}
The construction in
\texttt{size\_bounds/polynomial\_exceptional\_die.tex} uses $m=n-1$ old
dice and private node sets $Q_j$ for their density corrections. After
absorbing the factor $1/K$ into the signed weights, its common correction
functional is exactly of the form above. Any unused nodes only increase
the total distinguished-die size. Thus a nonzero correction forces
$K=\Omega(m^2)$ within this method.

If $B=0$, the uncorrected rational distinguished nodes already form a
scalar rational equal-weight interval quadrature of degree $m$.
The companion two-adic theorem then forces
$K\geq2^{\lfloor m/2\rfloor}$. In particular, asymptotically polynomial
versions of this disjoint-support construction necessarily use a nonzero
correction, and the quadratic limitation applies.

\paragraph{The remaining exceptional-die problem.}
Let $g(n)$ be the minimum number of faces of a prescribed die in an
$n$-die permutation-fair collection, with the other dice unrestricted.
The results obtained in this research session give
\[
 n-1\leq g(n)\leq Cn^2,
 \qquad g(2)=1,\quad g(3)=2,\quad g(4)=3.
\]
The proposition does not decide whether $g(n)$ has linear or quadratic
order. A construction with $o(n^2)$ distinguished faces would need a
mechanism beyond disjoint private supports representing one common
nonzero correction functional in this way.
\end{document}
